\begin{table}[!ht]
\centering
\caption{Escolha de hiperparâmetros e validação cruzada do XGBoost (Alencar, 2024;
Géron, 2021). Painel~A: seis configurações comparadas numa partição de validação
\emph{dentro} do treino --- o conjunto de teste permanece intocado. Painel~B: validação
cruzada $k$-\emph{fold} ($k=5$) no treino completo, com média e
desvio-padrão entre \emph{folds}. População completa: 6.155.273 observações de
treino e 1.538.819 de teste; nenhuma amostragem.}
\label{tab:ml_cv}
\small
\begin{tabular}{lccccc}
\multicolumn{6}{l}{\textbf{A. Busca na partição de validação}} \\
\toprule
Profundidade & Taxa de aprendizado & Árvores & $R^2$ & MAE & Sobreajuste \\
\midrule
\textbf{10} & 0,05 & 300 & 0,6281 & 0,4216 & 0,0109 \\
8 & 0,05 & 600 & 0,6253 & 0,4233 & 0,0063 \\
8 & 0,05 & 300 & 0,6210 & 0,4260 & 0,0034 \\
6 & 0,10 & 300 & 0,6186 & 0,4276 & 0,0015 \\
6 & 0,05 & 300 (atual) & 0,6145 & 0,4302 & 0,0007 \\
4 & 0,05 & 300 & 0,6037 & 0,4370 & −0,0001 \\
\bottomrule
\end{tabular}

\vspace{0.6em}
\begin{tabular}{lccc}
\multicolumn{4}{l}{\textbf{B. Validação cruzada no treino (5 \emph{folds})}} \\
\toprule
Configuração & $R^2$ (média $\pm$ dp) & MAE & RMSE \\
\midrule
Escolhida (profundidade 10) & 0,6279 $\pm$ 0,0006 & 0,4215 $\pm$ 0,0003 & 0,5632 $\pm$ 0,0004 \\
Anterior (profundidade 6) & 0,6142 $\pm$ 0,0006 & 0,4301 $\pm$ 0,0003 & 0,5735 $\pm$ 0,0005 \\
\bottomrule
\end{tabular}
\normalsize
\par\smallskip
\footnotesize\emph{Como ler:} ``Sobreajuste'' é a diferença entre o $R^2$ de treino e o de
validação --- valores próximos de zero indicam que o modelo não decorou os dados. A
configuração escolhida melhora o $R^2$ em 1,4 ponto percentual sem aumentar o
sobreajuste, e a variação entre \emph{folds} é da ordem de 0,0006 --- com 6,2~milhões de
observações de treino, o desempenho praticamente não depende de qual pedaço dos dados é
usado para validar. No conjunto de teste (intocado durante a escolha), $R^2 = 0,6280$.
Na unidade original, o erro mediano de previsão é de R\$~540 por mês
(37\% do rendimento observado), com correção de Duan para a
retransformação do logaritmo.
\end{table}
